%==============================================================
%  OPS260S - Group 1 project report  (master file)
%  Style: identical to the AD DS chapter (ADDS.tex)
%
%  HOW THIS PROJECT IS ORGANISED
%    main.tex            -> preamble, front matter, list of chapters
%    chapters/*.tex      -> one file per chapter (edit only your own)
%    images/             -> ALL screenshots live here
%    README.md           -> how to add images / new chapters
%
%  Compile with:  latexmk -pdf main.tex     (or just press Recompile in Overleaf)
%==============================================================
\documentclass[12pt,a4paper]{report}

%==============================================================
%  PACKAGES
%==============================================================
\usepackage[utf8]{inputenc}
\usepackage[T1]{fontenc}
\usepackage{mathpazo}                 % Palatino body text
\usepackage[scaled=0.95]{helvet}      % Helvetica for headings / captions
\usepackage{microtype}
\usepackage[a4paper,top=2.5cm,bottom=2.4cm,left=2.3cm,right=2.3cm,headheight=15pt]{geometry}
\usepackage{setspace}
\usepackage{xcolor}
\usepackage{graphicx}
\usepackage{float}                    % [H] keeps each screenshot next to its text
\usepackage{booktabs}
\usepackage{array}
\usepackage{enumitem}
\usepackage[font={small,sf},labelfont={bf,color=accent},labelsep=period,justification=centering,skip=6pt]{caption}
\usepackage[most]{tcolorbox}          % also loads listings
\usepackage{fancyhdr}
\usepackage{titlesec}
\usepackage[colorlinks=true,linkcolor=accent,citecolor=accent,urlcolor=accent,
            pdftitle={Design and Implementation of a Secure Windows Server 2019 Infrastructure for Enterprise Network Services},
            pdfauthor={Group 1}]{hyperref}
\usepackage{xurl}                     % loaded after hyperref so long URLs break

%==============================================================
%  COLOURS
%==============================================================
\definecolor{accent}{HTML}{1F4E79}
\definecolor{accentlight}{HTML}{EAF1F8}
\definecolor{warnred}{HTML}{B03A2E}
\definecolor{okgreen}{HTML}{1E7B45}

%==============================================================
%  PAGE & TEXT SETTINGS
%==============================================================
\setstretch{1.12}
\setlength{\parindent}{0pt}
\setlength{\parskip}{0.7em}
\emergencystretch=3em
\setlength{\fboxsep}{2pt}
\setlength{\fboxrule}{0.6pt}
\setlist{itemsep=2pt,topsep=4pt,parsep=0pt}
\renewcommand{\chapterautorefname}{Chapter}

%==============================================================
%  HEADINGS
%==============================================================
\titleformat{\chapter}[display]
  {\normalfont\sffamily\bfseries\color{accent}}
  {\Large\chaptertitlename\ \thechapter}
  {6pt}
  {\Huge}
  [\vspace{4pt}{\color{accent}\titlerule[1.4pt]}]
\titlespacing*{\chapter}{0pt}{-10pt}{22pt}

\titleformat{\section}
  {\normalfont\sffamily\Large\bfseries\color{accent}}
  {\thesection}{0.8em}{}
  [\vspace{-2pt}{\color{accent!35}\titlerule[0.5pt]}]
\titlespacing*{\section}{0pt}{20pt}{8pt}

\titleformat{\subsection}
  {\normalfont\sffamily\large\bfseries\color{accent!85!black}}
  {\thesubsection}{0.8em}{}
\titlespacing*{\subsection}{0pt}{16pt}{6pt}

%==============================================================
%  HEADER / FOOTER
%==============================================================
\renewcommand{\chaptermark}[1]{\markboth{\thechapter.\ #1}{}}
\pagestyle{fancy}
\fancyhf{}
\fancyhead[L]{\sffamily\small\color{gray}OPS260S Operating Systems 2}
\fancyhead[R]{\sffamily\small\color{gray}\leftmark}
\fancyfoot[C]{\sffamily\small\thepage}
\renewcommand{\headrulewidth}{0.4pt}
\renewcommand{\headrule}{\color{accent!40}\hrule width\headwidth height\headrulewidth}
\fancypagestyle{plain}{%
  \fancyhf{}\fancyfoot[C]{\sffamily\small\thepage}%
  \renewcommand{\headrulewidth}{0pt}}

%==============================================================
%  CUSTOM COMMANDS
%==============================================================

% ---- Screenshot placeholder box (shown automatically if the image is missing)
\newcommand{\placeholderbox}[1]{%
  \begin{tcolorbox}[enhanced,width=0.96\linewidth,center,height=5cm,valign=center,halign=center,
    colback=accentlight,frame hidden,boxrule=0pt,sharp corners,
    overlay={\draw[accent!60,dashed,line width=0.9pt] (frame.south west) rectangle (frame.north east);}]
    {\sffamily\bfseries\color{accent}SCREENSHOT PLACEHOLDER}\par\smallskip
    {\sffamily\small Save the image as}\par
    {\ttfamily\small images/\detokenize{#1}}\par\smallskip
    {\sffamily\footnotesize then recompile --- it will appear here automatically}
  \end{tcolorbox}}

% ---- Log of images that are still missing (written to missing_images.txt on every compile)
\newwrite\missingfile
\AtBeginDocument{%
  \immediate\openout\missingfile=missing_images.txt
  \immediate\write\missingfile{Screenshots still missing - save each file into the images/ folder:}%
  \immediate\write\missingfile{}}
\AtEndDocument{\immediate\closeout\missingfile}

% \shot{file}{caption}{label}  ->  numbered, captioned, framed, full-width screenshot.
%   If images/<file> exists it is inserted; otherwise a placeholder box is shown instead
%   and the file name is added to missing_images.txt.
\newcommand{\shot}[3]{%
  \begin{figure}[H]
    \centering
    \IfFileExists{images/#1}{%
      \fcolorbox{accent!50}{white}{\includegraphics[width=0.96\linewidth,height=0.5\textheight,keepaspectratio]{images/#1}}%
    }{%
      \placeholderbox{#1}%
      \immediate\write\missingfile{#1}%
    }%
    \caption{#2}\label{#3}
  \end{figure}}

% Blue side-bar call-out box:  \begin{callout}{Title} ... \end{callout}
\newtcolorbox{callout}[1]{enhanced,breakable,frame hidden,sharp corners,
  borderline west={3pt}{0pt}{accent},colback=accentlight,
  left=10pt,right=8pt,top=6pt,bottom=6pt,
  before upper={\textsf{\bfseries\color{accent}#1}\par\smallskip}}

% Troubleshooting entry:  \issue{title}{symptom}{fix}
\newcommand{\issue}[3]{%
  \begin{tcolorbox}[enhanced,breakable,colback=white,colframe=accent,
    colbacktitle=accent,coltitle=white,fonttitle=\sffamily\bfseries,
    title={#1},boxrule=0.6pt,arc=2pt,left=8pt,right=8pt,top=6pt,bottom=6pt]
    \textsf{\bfseries\color{warnred}Symptom.} #2\par\smallskip
    \textsf{\bfseries\color{okgreen}Fix.} #3
  \end{tcolorbox}}

% Reminder box for unfinished parts.  Switch all of them off with \showtodofalse
\newif\ifshowtodo \showtodotrue
\newcommand{\todo}[1]{%
  \ifshowtodo
  \begin{tcolorbox}[enhanced,breakable,colback=warnred!6,colframe=warnred,boxrule=0.5pt,
    arc=2pt,left=8pt,right=8pt,top=4pt,bottom=4pt]
    \textsf{\bfseries\color{warnred}TODO.} #1
  \end{tcolorbox}%
  \fi}

% PowerShell / command listings
\lstdefinestyle{ps}{basicstyle=\ttfamily\footnotesize,backgroundcolor=\color{accentlight},
  frame=single,rulecolor=\color{accent!30},framesep=5pt,xleftmargin=5pt,xrightmargin=5pt,
  breaklines=true,columns=fullflexible,keepspaces=true,showstringspaces=false,upquote=true,
  morecomment=[l]{\#},commentstyle=\color{okgreen}\itshape,aboveskip=0.6em,belowskip=0.4em}
\lstset{style=ps}

% Signature block for the declaration pages:  \sigblock{file}{Name}{Student no.}
\newcommand{\sigblock}[3]{%
  \begin{minipage}[t]{0.47\textwidth}\centering
    \IfFileExists{images/#1}{\includegraphics[height=1.5cm,width=0.9\linewidth,keepaspectratio]{images/#1}}%
      {\rule{0pt}{1.5cm}}\\[-1pt]
    \rule{\linewidth}{0.4pt}\\[-1pt]
    {\small #2, #3}
  \end{minipage}}

%==============================================================
%  DOCUMENT
%==============================================================
\begin{document}

%--------------------------------------------------------------
%  TITLE PAGE
%--------------------------------------------------------------
\begin{titlepage}
\centering\singlespacing
\vspace*{0cm}
\IfFileExists{images/cput_logo.png}{%
  \includegraphics[width=7.5cm]{images/cput_logo.png}}{%
  \begin{tcolorbox}[enhanced,width=7.5cm,center,halign=center,valign=center,height=2.2cm,
    colback=accentlight,frame hidden,boxrule=0pt,sharp corners,
    overlay={\draw[accent!60,dashed,line width=0.9pt] (frame.south west) rectangle (frame.north east);}]
    \sffamily\small Logo placeholder: save as \texttt{images/cput\_logo.png}
  \end{tcolorbox}}

\vspace{0.8cm}
{\sffamily\large\color{accent}\textls[80]{\MakeUppercase{Department of Electrical, Electronic\\ and Computer Engineering}}\par}

\vspace{0.9cm}
{\color{accent}\rule{0.78\textwidth}{1.4pt}}\par
\vspace{0.4cm}
{\sffamily\bfseries\huge\color{accent}Design and Implementation of a Secure Windows Server 2019 Infrastructure for Enterprise Network Services\par}
\vspace{0.4cm}
{\color{accent}\rule{0.78\textwidth}{1.4pt}}\par

\vspace{0.8cm}
{\large by}\par
\vspace{0.2cm}
{\sffamily\bfseries\LARGE Group 1\par}

\vspace{0.8cm}
Group project report submitted in partial fulfilment of the requirements for the degree\par
{\bfseries Bachelor of Engineering Technology in Computer Engineering}\par
in the Faculty of Engineering and the Built Environment\par
at the Cape Peninsula University of Technology\par

\vspace{0.5cm}
{\small
\begin{tabular}{@{}r@{\hspace{0.8em}}l@{}}
  Examiner: & Mr.\ Vusumuzi Moyo \\
  Course Code: & OPS260S -- OPERATING SYSTEMS 2 \\
\end{tabular}\par}

\vspace{0.4cm}
Bellville Campus\par
September 2026\par

\vfill
{\footnotesize The report may not be published either in part (in scholarly, scientific or technical journals), or as a whole (as a monograph), unless permission has been obtained from the University.\par}
\end{titlepage}

%--------------------------------------------------------------
%  FRONT MATTER  (roman page numbers)
%--------------------------------------------------------------
\pagenumbering{roman}

\chapter*{Declaration of Authorship}
We, the undersigned, declare that the contents of this report represent our own unaided work, except where explicitly acknowledged. The report has not previously been submitted for academic examination towards any qualification. Furthermore, it represents our own opinions and not necessarily those of the Cape Peninsula University of Technology.

\vspace{1.2cm}
\sigblock{sig_nguyen.png}{Anh Duy Nguyen}{251163733}\hfill
\sigblock{sig_lamprecht.png}{Carl Lamprecht}{205133681}\\[1.2em]
\sigblock{sig_ngesi.png}{Peter Hlalubonga Ngesi}{250281228}\hfill
\sigblock{sig_koekemoer.png}{Gian Koekemoer}{250060248}\\[1.2em]
\sigblock{sig_koch.png}{Tihmo Koch}{250030993}\hfill
\sigblock{sig_kheswa.png}{Sibusiso Isaac Kheswa}{250175282}

\vspace{1.4cm}
October 26, 2026\\[-2pt]
\rule{0.47\textwidth}{0.4pt}\\[-2pt]
{\small Date}

\clearpage
\chapter*{AI Declaration}
We declare that generative artificial intelligence (AI) tools were used in the preparation of this report as follows:
\begin{itemize}
  \item \textbf{Grammar, spelling, and language refinement} (e.g., rephrasing for clarity and readability).
  \item \textbf{\LaTeX{} formatting assistance} (structure and code improvements).
\end{itemize}
All AI-generated suggestions were carefully reviewed, edited, and verified by the authors. The core ideas, analysis, research, and final content are entirely the work of the authors. No AI tool was used to generate original ideas, conduct the literature review, or produce the main technical content without substantial human input and oversight.

This declaration is made in accordance with academic integrity principles and the Cape Peninsula University of Technology guidelines on responsible AI use.

\vspace{0.8cm}
\sigblock{sig_nguyen.png}{Anh Duy Nguyen}{251163733}\hfill
\sigblock{sig_lamprecht.png}{Carl Lamprecht}{205133681}\\[1.2em]
\sigblock{sig_ngesi.png}{Peter Hlalubonga Ngesi}{250281228}\hfill
\sigblock{sig_koekemoer.png}{Gian Koekemoer}{250060248}\\[1.2em]
\sigblock{sig_koch.png}{Tihmo Koch}{250030993}\hfill
\sigblock{sig_kheswa.png}{Sibusiso Isaac Kheswa}{250175282}

\vspace{1.2cm}
October 26, 2026\\[-2pt]
\rule{0.47\textwidth}{0.4pt}\\[-2pt]
{\small Date}

\clearpage
\chapter*{Abstract}
This report presents the installation, configuration, and testing of a Windows Server 2019 environment developed to meet the network and security requirements of CPUT. The project was implemented using the hypervisor utilised during Semester~1 and included the deployment of several key server roles and services.

The implementation involved configuring Active Directory Domain Services (AD DS) for centralized user and resource management, Domain Name System (DNS) for name resolution, and Dynamic Host Configuration Protocol (DHCP) for automatic IP address allocation. In addition, Group Policy settings were applied to improve the security of the environment by restricting Command Prompt access, disabling removable storage devices, preventing unauthorized software installations, and enforcing stronger password policies.

Remote Access Services and Windows Certificate Services were also configured to provide secure connectivity and authentication within the network environment. Each service was tested to verify that it functioned correctly and met the specified requirements. Screenshots of all installation, configuration, and testing procedures were captured and included as evidence of successful implementation.

The results demonstrate how Windows Server 2019 can be used to create a secure and manageable network infrastructure that supports organisational requirements while following recommended security practices.

\vspace{0.6em}
\textsf{\bfseries\color{accent}Keywords:} CPUT, Windows Server 2019, Active Directory, DNS, DHCP, Group Policy, Remote Access, Certificate Services, Network Administration, Security Policies, Server Configuration

\clearpage
\tableofcontents
\clearpage
\listoffigures
\clearpage
\listoftables

%--------------------------------------------------------------
%  MAIN MATTER  (arabic page numbers)
%--------------------------------------------------------------
\clearpage
\pagenumbering{arabic}

\chapter{Introduction}
\label{chap:intro}

As organisations continue to grow, the need for a reliable and secure network infrastructure becomes increasingly important. CPUT requires an IT environment that can support its expanding operations while ensuring efficient management of users, devices, and network resources. To meet these requirements, a Windows Server 2019 environment was implemented using the hypervisor used during Semester~1.

This project focuses on the installation, configuration, and testing of key Windows Server services, including Active Directory Domain Services (AD DS), DNS, DHCP, Group Policy, Remote Access, and Windows Certificate Services. Screenshots were captured throughout the process to provide evidence of the installation and configuration steps performed.

\section{Background}
\label{sec:intro-background}
Windows Server 2019 is widely used in organisations to manage network resources, user accounts, security policies, and shared services from a central location. Services such as Active Directory, DNS, and DHCP help administrators manage users and devices more efficiently while ensuring that network communication functions correctly.

Security is also an important part of network administration. For this reason, several security policies were implemented in this project. These include disabling Command Prompt access, preventing the use of removable storage devices, restricting software installations, enforcing a minimum password length of ten characters, and requiring users to change their passwords regularly. These measures help reduce security risks and improve the overall protection of the network environment.

\section{Objectives}
\label{sec:intro-objectives}
The objectives of this report are:
\begin{itemize}
  \item To install and configure Windows Server 2019.
  \item To configure and test Active Directory Domain Services (AD DS).
  \item To configure and test DHCP services.
  \item To configure and test DNS services.
  \item To implement Group Policy settings that improve system security.
  \item To configure Remote Access services.
  \item To install and configure Windows Certificate Services.
  \item To verify that all configured services are functioning correctly.
  \item To document all installation, configuration, and testing procedures using screenshots.
\end{itemize}

\section{Report Structure}
\label{sec:intro-structure}
The remainder of this report is organised as follows:
\begin{itemize}
  \item Chapter~\ref{chap:install} covers the installation and initial configuration of Windows Server 2019. % Duy
  \item Chapter~\ref{chap:adds} discusses the configuration and testing of Active Directory Domain Services (AD DS). % Gian
  \item Chapter~\ref{chap:dhcp} explains the configuration and testing of DHCP services. % Peter
  \item Chapter~\ref{chap:dns} presents the configuration and testing of DNS services. % tbd
  \item Chapter~\ref{chap:gpo} describes the implementation of Group Policy settings. % Carl
  \item Chapter~\ref{chap:remote} discusses the configuration of Remote Access services. % Isaac
  \item Chapter~\ref{chap:certs} explains the installation and configuration of Windows Certificate Services. % Tihmo
  \item Chapter~\ref{chap:conclusion} provides the conclusion of the project. % Duy
  \item The final sections contain the bibliography and the appendix.
\end{itemize}

\chapter{Installation and Initial Configuration of Windows Server 2019}
\label{chap:install}

\section{Introduction}
\label{sec:install-intro}
This chapter describes the installation and initial configuration of Windows Server 2019 within the virtualised environment used during Semester~1. The purpose of this stage was to establish a functioning server that could support the deployment of the network services required for the project.

The installation process involved preparing the virtual machine using VirtualBox, installing the operating system, and completing the initial server configuration. These steps provided the foundation for the subsequent configuration of Active Directory Domain Services (AD DS), DNS, DHCP, Group Policy, Remote Access, and Windows Certificate Services. \autoref{tab:install-config} lists the values used throughout this chapter.

\begin{table}[H]
  \centering
  \small\sffamily
  \renewcommand{\arraystretch}{1.3}
  \begin{tabular}{@{}p{5.6cm}p{8.6cm}@{}}
    \toprule
    \textbf{Setting} & \textbf{Value used} \\
    \midrule
    Hypervisor & Oracle VirtualBox \\
    Virtual machine name & WinServer \\
    Guest operating system & Windows Server 2019 (64-bit) \\
    Edition installed & Standard Evaluation (Desktop Experience) \\
    Base memory & 6144 MB \\
    Processors & 3 \\
    Virtual hard disk & 32 GB, VDI (dynamically allocated) \\
    Network adapter & Intel PRO/1000 MT Desktop (NAT) \\
    Installer language / time and currency / keyboard & English (United States) / English (South Africa) / US \\
    \bottomrule
  \end{tabular}
  \caption{Key configuration values used in this chapter}\label{tab:install-config}
\end{table}

\section{Virtual Machine Preparation}
\label{sec:install-vm}
Before installing Windows Server 2019, a virtual machine was prepared using VirtualBox. Virtualisation allows multiple operating systems to run on a single physical computer while maintaining separate virtual environments. The New Virtual Machine wizard was opened (\autoref{fig:install-01}), and the machine was given a name and pointed at the Windows Server 2019 installation image (\autoref{fig:install-02}).

\shot{install_01.png}{Initial Virtual Machine Configuration}{fig:install-01}
\shot{install_02.png}{VM Name and Windows Server 2019 ISO Selected}{fig:install-02}

The virtual machine was configured with the necessary resources to support the operating system and the planned server roles. These resources included virtual memory and processor allocation (\autoref{fig:install-03}), virtual storage (\autoref{fig:install-04}), and a network adapter. The Windows Server 2019 installation media was then attached to the virtual machine, allowing it to boot from the installation image. The finished machine is listed in the VirtualBox Manager (\autoref{fig:install-05}), and its final settings are shown in \autoref{fig:install-06}.

\shot{install_03.png}{Virtual Hardware: Memory and CPUs}{fig:install-03}
\shot{install_04.png}{Virtual Hard Disk Configuration}{fig:install-04}
\shot{install_05.png}{WinServer Listed in the VirtualBox Manager}{fig:install-05}
\shot{install_06.png}{Final Virtual Machine Configuration Summary}{fig:install-06}

\section{Installation of Windows Server 2019}
\label{sec:install-os}
Once the virtual machine had been configured, it was started and the Windows Server 2019 installation process was initiated. The installation wizard was used to select the appropriate language, time and currency format, and keyboard layout (\autoref{fig:install-lang}), after which the installation was started (\autoref{fig:install-07}).

\shot{install_14.png}{Language, Time and Currency Format, and Keyboard Selection}{fig:install-lang}
\shot{install_07.png}{Windows Setup: Install Now}{fig:install-07}

The required Windows Server 2019 edition was then selected (\autoref{fig:install-08}). The installation proceeded by accepting the applicable licence terms (\autoref{fig:install-09}), selecting the installation type (\autoref{fig:install-10}), and choosing the virtual disk on which the operating system would be installed (\autoref{fig:install-11}).

\shot{install_08.png}{Selecting the Operating System Edition}{fig:install-08}
\shot{install_09.png}{Accepting the Licence Terms}{fig:install-09}
\shot{install_10.png}{Selecting the Installation Type}{fig:install-10}
\shot{install_11.png}{Selecting the Installation Disk}{fig:install-11}

After the installation files had been copied and the required setup operations completed (\autoref{fig:install-12}), the virtual machine restarted (\autoref{fig:install-13}). The initial administrator password was configured (\autoref{fig:install-15}), allowing access to the server using the local Administrator account (\autoref{fig:install-16}).

\shot{install_12.png}{Installing Windows}{fig:install-12}
\shot{install_13.png}{Restart Prompt After Installation}{fig:install-13}
\shot{install_15.png}{Setting the Administrator Password}{fig:install-15}
\shot{install_16.png}{Administrator Log-In Screen}{fig:install-16}
\shot{install_17.png}{Windows Server 2019 Desktop After First Log-In}{fig:install-17}

\section{Testing}
\label{sec:install-testing}
Testing was conducted to verify that Windows Server 2019 had been installed successfully and that the initial server configuration was functioning as expected. The tests focused on operating system availability, server identification, network configuration, and basic network connectivity.

\subsection{Operating System Verification}
\label{sec:install-osverify}
The server was restarted to confirm that Windows Server 2019 could boot successfully after installation. The administrator was able to log in, and the Server Manager dashboard was accessible (\autoref{fig:install-18}). This confirmed that the operating system was operational and ready for further configuration.

\shot{install_18.png}{Server Manager Dashboard After the Installation of Windows Server 2019}{fig:install-18}

The dashboard provided access to the server's management tools and the interface used to install and configure additional server roles and features.

Further testing should verify the server hostname, IP configuration, and network connectivity. The following commands can be used to perform these checks:
\begin{itemize}
  \item \texttt{hostname} -- displays the configured computer name.
  \item \texttt{ipconfig /all} -- displays the network adapter configuration.
  \item \texttt{ping <gateway-IP>} -- tests connectivity to the configured network gateway, where applicable.
\end{itemize}

The outputs should be recorded as supporting evidence. Any failed tests should be investigated and resolved before proceeding with the configuration of additional server roles.

More comprehensive testing of Active Directory Domain Services, DNS, DHCP, Remote Access, and Windows Certificate Services is addressed in the relevant chapters of this report.

\section{Chapter Summary}
\label{sec:install-summary}
This chapter described the preparation of the virtual machine, installation of Windows Server 2019, and completion of the initial server configuration. The process established the operating system environment required for the deployment of additional server roles and network services.

Initial verification procedures were also described to assess operating system availability, server identification, network configuration, and connectivity. These steps establish the foundation for the remaining stages of the project.

The next chapter discusses the configuration and testing of Active Directory Domain Services (AD DS), which provides centralised management of users, computers, and authentication within the network environment.

\chapter{Active Directory Domain Services (AD DS) Configuration}
\label{chap:adds}

\section{What is AD DS?}
\label{sec:what}
Active Directory Domain Services (AD DS) is the foundational directory service for Microsoft Active Directory. It operates as a centralized identity management system that authenticates users, computers, and devices on a Windows network. It organizes network resources hierarchically and verifies credentials, ensuring that only authorized entities can access network assets \cite{msft-adds-overview}.

\section{Why is AD DS important?}
\label{sec:why}
For a growing infrastructure like CPUT's, AD DS is critical because it centralizes network administration. Instead of managing user accounts and security policies on every individual machine, administrators can manage the entire network from a single domain controller. Crucially, AD DS forms the backbone required to deploy Group Policy Objects (GPOs). Without AD DS, enforcing the required project security policies---such as disabling the Command Prompt, blocking removable media, preventing unauthorized software installations, and enforcing strict password length (10+ characters) and age rules (under 42 days)---would be virtually impossible to scale effectively across the organization \cite{msft-gpo-sec}.

\section{Security Risks if Not Set Up or Set Up Incorrectly}
\label{sec:risks}
If AD DS is omitted, the network remains a disjointed workgroup with no centralized security, leaving endpoints highly vulnerable to local exploitation, malware from USB drives, and unauthorized software. If AD DS is set up incorrectly (e.g., weak administrator credentials, misconfigured DNS, or overly permissive access rights), it becomes a single point of failure. An attacker who compromises a misconfigured Domain Controller gains ``keys to the kingdom,'' allowing them to silently escalate privileges, distribute malware laterally across all connected machines, and bypass all intended security policies \cite{msft-sec-guidelines}.

\section{Step-by-Step Configuration Guide}
\label{sec:guide}
The configuration is carried out in four steps: promoting the server (\autoref{sec:step1}), verifying the installation (\autoref{sec:step2}), creating a reverse lookup zone (\autoref{sec:step3}), and testing DNS resolution (\autoref{sec:step4}). \autoref{tab:config} lists the values used throughout this guide. If anything goes wrong, see \autoref{sec:troubleshooting}.

\begin{table}[H]
  \centering
  \small\sffamily
  \renewcommand{\arraystretch}{1.3}
  \begin{tabular}{@{}p{5.6cm}p{8.6cm}@{}}
    \toprule
    \textbf{Setting} & \textbf{Value used} \\
    \midrule
    Root domain name & OPS.local \\
    NetBIOS domain name & OPS \\
    Forest / domain functional level & Windows Server 2016 \\
    Server static IP (example) & 192.168.0.5 \\
    Reverse lookup network ID (example) & 192.168.0 \\
    Reverse zone type & Primary, stored in Active Directory \\
    Reverse zone replication & All DNS servers on domain controllers in this domain \\
    Dynamic updates & Secure only \\
    \bottomrule
  \end{tabular}
  \caption{Key configuration values used in this chapter}\label{tab:config}
\end{table}

%--------------------------------------------------------------
\subsection{Step 1: Promote the Server to a Domain Controller}
\label{sec:step1}
Begin the AD DS Configuration Wizard by selecting ``Add a new forest'' (\autoref{fig:forest}) and specifying your Root domain name, e.g., OPS.local (\autoref{fig:rootdomain}).

\shot{ADDS_01.png}{Deployment Configuration - Adding a new forest}{fig:forest}
\shot{ADDS_02.png}{Specifying the Root Domain Name}{fig:rootdomain}

Set the Forest and Domain functional levels to Windows Server 2016. Ensure DNS server and Global Catalog are checked, then enter and confirm a strong Directory Services Restore Mode (DSRM) password (\autoref{fig:dsrm}).

\shot{ADDS_03.png}{Domain Controller Options and DSRM Password}{fig:dsrm}

\begin{callout}{Keep the DSRM password safe}
The DSRM password is required to repair or restore the directory later. Record it in a secure location and do not reuse another account's password.
\end{callout}

You will encounter a warning stating a delegation for this DNS server cannot be created (\autoref{fig:delegation}). This is expected for the first DC in a new forest (see \autoref{sec:troubleshooting}, item~2). Click Next.

\shot{ADDS_04.png}{DNS Delegation Warning}{fig:delegation}

Verify the assigned NetBIOS domain name, e.g., OPS (\autoref{fig:netbios}).

\shot{ADDS_05.png}{NetBIOS Domain Name Verification}{fig:netbios}

Review your selected AD DS configuration settings before proceeding (\autoref{fig:review}). You can click ``View script'' to see the underlying PowerShell commands (\autoref{fig:ps1}, \autoref{fig:review2} and \autoref{fig:ps2}).

\shot{ADDS_06.png}{Review Options Summary}{fig:review}
\shot{ADDS_07.png}{PowerShell Deployment Script - Part 1}{fig:ps1}
\shot{ADDS_08.png}{Review Options - Continued}{fig:review2}
\shot{ADDS_09.png}{PowerShell Deployment Script - Part 2}{fig:ps2}

The system will verify prerequisites. You may see warnings regarding default cryptography algorithms and physical network adapters without static IPs (\autoref{fig:prereqwarn}), but ensure you see the green checkmark stating ``All prerequisite checks passed successfully'' (\autoref{fig:prereqok}). Click Install. If the static IP warning persists, see \autoref{sec:troubleshooting}, item~1.

\shot{ADDS_10.png}{Prerequisites Check Warning Details}{fig:prereqwarn}
\shot{ADDS_11.png}{Prerequisites Check Passed}{fig:prereqok}

The server will begin configuring the DNS Server service and AD DS, automatically rebooting once complete (\autoref{fig:install}).

\shot{ADDS_12.png}{AD DS Installation Progress}{fig:install}

%--------------------------------------------------------------
\subsection{Step 2: Verify Active Directory Installation}
\label{sec:step2}
After reboot, log in as the Domain Administrator. Open Server Manager, click ``Tools,'' and select ``Active Directory Users and Computers'' (\autoref{fig:aduc}).

\shot{ADDS(3).png}{Accessing AD Users and Computers}{fig:aduc}

Expand the tree to verify your domain (OPS.local) is present (\autoref{fig:domaincontainer}), then right-click your domain name and select ``Properties'' (\autoref{fig:domainprops}).

\shot{ADDS(4).png}{Domain Container Verification}{fig:domaincontainer}
\shot{ADDS(5).png}{Opening Domain Properties}{fig:domainprops}

Check the General tab to confirm the Domain and Forest functional levels are active (\autoref{fig:funclevel}).

\shot{ADDS(6).png}{Domain Functional Level Verification}{fig:funclevel}

%--------------------------------------------------------------
\subsection{Step 3: Configure DNS and Create a Reverse Lookup Zone}
\label{sec:step3}
Return to Server Manager, click ``Tools,'' and select ``DNS'' (\autoref{fig:dnsopen}). Expand your server name and verify that the Forward Lookup Zone for your domain was created automatically (\autoref{fig:fwdzones}).

\shot{ADDS(7).png}{Opening DNS Manager}{fig:dnsopen}
\shot{ADDS(8).png}{Forward Lookup Zones}{fig:fwdzones}

Right-click ``Reverse Lookup Zones'' and select ``New Zone...'' (\autoref{fig:newrev} and \autoref{fig:zonewelcome}).

\shot{ADDS(9).png}{Creating a New Reverse Lookup Zone}{fig:newrev}
\shot{ADDS(10).png}{New Zone Wizard Welcome Screen}{fig:zonewelcome}

Select ``Primary zone'' and ensure ``Store the zone in Active Directory'' is checked (\autoref{fig:zonetype}) \cite{msft-dns-integration}.

\shot{ADDS(11).png}{Zone Type Selection}{fig:zonetype}

Set replication to ``To all DNS servers running on domain controllers in this domain'' (\autoref{fig:zonescope}).

\shot{ADDS(12).png}{Active Directory Zone Replication Scope}{fig:zonescope}

Choose ``IPv4 Reverse Lookup Zone'' (\autoref{fig:ipv4rev}) and enter your network ID, e.g., 192.168.0 (\autoref{fig:netid}).

\shot{ADDS(13).png}{IPv4 Reverse Lookup Selection}{fig:ipv4rev}
\shot{ADDS(14).png}{Network ID Configuration}{fig:netid}

Select ``Allow only secure dynamic updates'' to comply with AD security best practices (\autoref{fig:dynupdates}), then click Finish (\autoref{fig:zonefinish}). The new zone is confirmed in \autoref{fig:revverify}.

\shot{ADDS(15).png}{Dynamic Update Security Configuration}{fig:dynupdates}
\shot{ADDS(16).png}{Completing the New Zone Wizard}{fig:zonefinish}
\shot{ADDS(17).png}{Reverse Lookup Zone Verification}{fig:revverify}

%--------------------------------------------------------------
\subsection{Step 4: Update DNS Records and Test Resolution}
\label{sec:step4}
Navigate to your Forward Lookup Zone, locate your server's Host (A) record, right-click, and select Properties (\autoref{fig:arecord}). Check the box that says ``Update associated pointer (PTR) record'' and click OK (\autoref{fig:ptrenable}).

\shot{ADDS(18).png}{Host (A) Record Properties}{fig:arecord}
\shot{ADDS(19).png}{Enabling PTR Record Update}{fig:ptrenable}

Refresh the Reverse Lookup Zone to confirm the Pointer (PTR) record has successfully populated (\autoref{fig:ptrverify}).

\shot{ADDS(20).png}{PTR Record Verification}{fig:ptrverify}

Open your Network Connections, access the IPv4 properties of your adapter, and set the Preferred DNS server to your server's static IP address, e.g., 192.168.0.5 (\autoref{fig:adapterdns}).

\shot{ADDS(21).png}{Network Adapter DNS Configuration}{fig:adapterdns}

Open PowerShell as Administrator. Type \texttt{nslookup} followed by your server's fully qualified domain name (FQDN) to test forward and reverse lookups (\autoref{fig:nslookupbasic} and \autoref{fig:nslookupfull}). Successful replies confirm DNS and AD DS are communicating correctly.

\shot{ADDS(22).png}{Basic DNS Resolution Test}{fig:nslookupbasic}
\shot{ADDS(23).png}{Comprehensive nslookup Verification}{fig:nslookupfull}

%--------------------------------------------------------------
\section{Troubleshooting: Possible Sticking Points \& Fixes}
\label{sec:troubleshooting}

\issue{1. Static IP Prerequisites Warning or Failure}
{The Prerequisites Check throws a hard error or persistent warning regarding physical network adapters not having static IP addresses.}
{Before promoting the server, open ``Network and Sharing Center'' $\rightarrow$ ``Change adapter settings.'' Right-click your network interface, select Properties, and manually assign an IPv4 address, Subnet Mask, Default Gateway, and set the Preferred DNS to 127.0.0.1 (loopback) or the static IP you just assigned.}

\issue{2. DNS Delegation Warning During AD DS Setup}
{A yellow warning icon appears on the DNS Options page stating ``A delegation for this DNS server cannot be created'' (\autoref{fig:delegation}).}
{This is entirely normal when building the first Domain Controller in a new forest, as there is no existing parent DNS infrastructure to integrate with. No action is required---simply click Next.}

\issue{3. nslookup Times Out or Fails to Resolve}
{When running \texttt{nslookup} in PowerShell, it returns ``Default Server: UnKnown'' or times out without resolving the IP/Hostname.}
{This usually means the server's network adapter is looking at an external DNS instead of itself. Go to your IPv4 Network Properties and ensure the ``Preferred DNS server'' is set to the server's own static IP, e.g., 192.168.0.5 (\autoref{fig:adapterdns}). Additionally, check the DNS Manager to ensure the PTR record was actually created in the Reverse Lookup Zone (\autoref{fig:ptrverify}). If missing, manually add it.}

\issue{4. Group Policy Restrictions Not Applying to Users}
{You configure GPOs to disable CMD, prevent software installations, and restrict removable media, but users can still bypass them.}
{Ensure that your test users/computers are actually moved into the appropriate Organizational Units (OUs) within Active Directory Users and Computers (\autoref{fig:aduc}). GPOs linked to custom OUs will not apply if the users remain in the default ``Users'' container. Run \texttt{gpupdate /force} in PowerShell on the client machine to apply policies immediately.}

\chapter{DHCP Configuration}
\label{chap:dhcp}

\section{Background}
\label{sec:dhcp-background}
Every machine, from computers to servers, requires an IP address to be able to send and receive packets of data \cite{krause2019}. These IP addresses act like home addresses on a street, and whilst DNS is responsible for resolving those addresses into names, there are two ways of assigning IP addresses to machines: static addressing and DHCP \cite{krause2019}.

Static addressing is the process of configuring IP addresses on a system manually. It is a hands-on approach where the address is entered directly into the NIC settings of a particular device \cite{krause2019}. Despite its simplicity, this method does not scale in environments that host many devices, such as a small or large company with laptops, desktops, printers, servers, and other machines. Each of these devices needs an IP address, but assigning them manually becomes tedious at scale. The solution to this problem is the Dynamic Host Configuration Protocol (DHCP) \cite{krause2019}.

DHCP is a protocol that provides IP addressing information to a large number of devices automatically: a machine that is plugged into the network can obtain its address information from the DHCP server \cite{krause2019}. When a laptop or phone connects to a Wi-Fi router, a DHCP server assigns it the addressing information it needs to route traffic on the network. On public Wi-Fi the DHCP server runs on the router itself, whereas in businesses where Windows Server acts as the datacentre, DHCP services are most often hosted on one or more servers across the network \cite{krause2019}.

\section{Configuration Steps}
\label{sec:dhcp-steps}
The configuration is carried out in two stages: the post-install configuration of the DHCP role (\autoref{sec:dhcp-post}) and the creation of a DHCP scope (\autoref{sec:dhcp-scope}). \autoref{tab:dhcp-config} lists the values used throughout this chapter. If the DHCP console cannot be used, or something goes wrong, see \autoref{sec:dhcp-issues}.

\begin{table}[H]
  \centering
  \small\sffamily
  \renewcommand{\arraystretch}{1.3}
  \begin{tabular}{@{}p{5.6cm}p{8.6cm}@{}}
    \toprule
    \textbf{Setting} & \textbf{Value used} \\
    \midrule
    Scope name & GROUP 1-SCOPE \\
    Address pool & 192.168.0.100 -- 192.168.0.254 \\
    Subnet mask & 255.255.255.0 (/24) \\
    Exclusions and delay & None \\
    Lease duration & 8 days (default) \\
    Default gateway (router) & 192.168.0.1 \\
    Parent domain & OPS.local \\
    DNS servers & 192.168.0.5 and 8.8.8.8 \\
    WINS servers & None \\
    Scope state & Activated \\
    Authorisation credentials & \texttt{OPS\textbackslash Administrator} \\
    \bottomrule
  \end{tabular}
  \caption{Key configuration values used in this chapter}\label{tab:dhcp-config}
\end{table}

%--------------------------------------------------------------
\subsection{DHCP Post-Install Configuration}
\label{sec:dhcp-post}
Like other server roles, DHCP must first be added to the server during the installation phase. After the DHCP role has been added, it is not yet properly configured.

The DHCP server can give clients two kinds of access groups: DHCP Administrators and DHCP Users \cite{dhcp-groups}. DHCP Users have read-only access to the server, which allows them to view, but not modify, server data, including the DHCP server configuration, registry keys, DHCP log files, and the database. They cannot create scopes, modify option values, or create reservations or exclusion ranges. DHCP Administrators are able to view and modify data at the DHCP server. They can create and delete scopes, add reservations, change option values, create superscopes, or perform any other activity needed to administer the DHCP server. Despite this control, they do not have unlimited administrative rights, because they only have rights on the local computer \cite{dhcp-groups}.

These two groups are local to the server, which means they cannot authorise DHCP servers; this can only be done by Domain Admins \cite{dhcp-groups}. Authorising the server is needed because DHCP communication typically occurs before any user or computer authentication, and because DHCP is based on IP, an unknown DHCP server could hand invalid information to clients. Authorisation prevents this \cite{msnoob-dhcp}.

The authorisation step in the wizard requires the credentials of an account that is allowed to authorise the server in Active Directory (\autoref{fig:dhcp-01} and \autoref{fig:dhcp-02}). So far the only credential set was \texttt{OPS\textbackslash Administrator}, and any user who attempts to connect to the server in future must use this username. The wizard confirms the result in \autoref{fig:dhcp-03}.

\shot{dhcp_01.png}{Description of the DHCP Post-Install Configuration Wizard}{fig:dhcp-01}
\shot{dhcp_02.png}{Authorisation of the DHCP Post-Install Configuration Wizard}{fig:dhcp-02}
\shot{dhcp_03.png}{Post-Install Configuration Completed}{fig:dhcp-03}

\begin{callout}{Restart the DHCP service}
The wizard asks for the DHCP server service to be restarted so that the new security groups become effective. Do this before creating the scope.
\end{callout}

%--------------------------------------------------------------
\subsection{The DHCP Scope}
\label{sec:dhcp-scope}
For a DHCP server to hand out IP addresses, it needs to know which IP addresses, subnet mask, default gateway and DNS server addresses are valid within the network, so that it can package that information and hand it out to the computers that request it. This package of information is called a DHCP scope \cite{krause2019}.

To begin defining the scope, the DHCP management tool was launched from the Tools menu in Server Manager (\autoref{fig:dhcp-04}). This opens a window for DHCP with the name of the device that the server is running on (\autoref{fig:dhcp-05}).

\shot{dhcp_04.png}{DHCP Selected on the Tools Dropdown Menu}{fig:dhcp-04}
\shot{dhcp_05.png}{DHCP Window}{fig:dhcp-05}

After opening the DHCP window, the server was expanded to show two protocols: IPv4 and IPv6. IPv4 was selected, which displayed information about adding a scope, along with the server options, policies and filters. IPv4 did not have a scope yet, so one needed to be configured. IPv4 was right-clicked to open the menu containing the option to add a new scope (\autoref{fig:dhcp-06}). Choosing this option opened the New Scope Wizard (\autoref{fig:dhcp-07}). After clicking Next, the basics of the scope were decided. The name given was GROUP 1-SCOPE (\autoref{fig:dhcp-08}).

\shot{dhcp_06.png}{Adding a Scope}{fig:dhcp-06}
\shot{dhcp_07.png}{New Scope Wizard}{fig:dhcp-07}
\shot{dhcp_08.png}{Group One Scope Made}{fig:dhcp-08}

There exist 4.3 billion possible IPv4 addresses, of which roughly 3.7 billion are available for use \cite{ipme-ipv4}. A scope does not need to accommodate billions of addresses; it just needs a few hundred or a few dozen. When configuring it, the range of IP addresses had to be decided. The addresses that the scope will hand out are 192.168.0.100 to 192.168.0.254 (\autoref{fig:dhcp-10}), and this is the address pool. No exclusions or delays were made in the scope (\autoref{fig:dhcp-09}).

Alongside the handed-out IP addresses, the scope also has the option to configure the amount of time that a client keeps an assigned IP address before renewal, ensuring efficient management of IP addresses \cite{rapidseedbox-lease}. This period of time is called the lease duration. It was left at the default of 8 days (\autoref{fig:dhcp-11}).

\shot{dhcp_09.png}{No Exclusions Made}{fig:dhcp-09}
\shot{dhcp_10.png}{Setting the IP Address Range}{fig:dhcp-10}
\shot{dhcp_11.png}{Lease Duration of 8 Days Set}{fig:dhcp-11}

When clients use the server, access to the internet is an important feature. In order for the clients to be able to access the internet, they need to be able to talk to devices and networks outside of the DHCP subnet. A default gateway allows for this; without it, a client can only reach devices and networks inside its own subnet. The default gateway was given the IP address 192.168.0.1 (\autoref{fig:dhcp-12} and \autoref{fig:dhcp-13}).

\shot{dhcp_12.png}{Default Gateway Configured (1)}{fig:dhcp-12}
\shot{dhcp_13.png}{Default Gateway Configured (2)}{fig:dhcp-13}

After the default gateway, the next step was setting up the DNS server. The DHCP server hands out IP addresses and more, so that a client can use the network. An IP address allows a device to communicate with other hosts, but IP addresses are inconvenient to memorise or type. The DNS server and domain name exist so that devices can reach hosts by names that resolve to IP addresses \cite{msft-dhcp-dns}. The domain name was set to OPS.local, and two DNS server addresses, 192.168.0.5 and 8.8.8.8, were configured (\autoref{fig:dhcp-14}). This means that, in order to access the server, a client simply needs to type in the domain name, which makes it more accessible for non-domain-joined machines.

\shot{dhcp_14.png}{DNS Server and Domain Name}{fig:dhcp-14}

After the DNS server and domain name were configured, there was an option for WINS server configuration. WINS is a part of the scope that resolves NetBIOS computer names to IP addresses. It is needed to support pre-Windows~2000 systems and older applications that use NetBIOS over TCP/IP \cite{msft-wins}. None of the systems planned to be connected to the network qualify as pre-Windows~2000 systems, therefore WINS is not needed and was left vacant (\autoref{fig:dhcp-15}).

\shot{dhcp_15.png}{WINS Server Configuration Left Vacant}{fig:dhcp-15}

All necessary parts of the scope have now been configured. The only step left was the activation of the scope (\autoref{fig:dhcp-16}). Upon activating the scope, the wizard presents a completion display, and the Finish button was clicked to close it (\autoref{fig:dhcp-17}). The scope is now activated and ready for testing. It will hand out IP addresses from 192.168.0.100 to 192.168.0.254. At the time of creation there were no address leases yet (\autoref{fig:dhcp-18}).

\shot{dhcp_16.png}{Activation of the Scope}{fig:dhcp-16}
\shot{dhcp_17.png}{Finishing the New Scope Wizard}{fig:dhcp-17}
\shot{dhcp_18.png}{New DHCP Scope Created}{fig:dhcp-18}

%--------------------------------------------------------------
\section{Issues Encountered for DHCP Server Setup and Alternative Paths}
\label{sec:dhcp-issues}

\issue{DHCP Console Stops Functioning}
{The DHCP console can stop working, for example after a domain name change or after a new Ethernet adapter is added for internet access, so a new scope cannot be configured through the window.}
{Windows PowerShell is still available on Windows Server 2019 and bypasses surface-level console problems. The same configuration can be done with PowerShell commands (see \autoref{sec:dhcp-ps} and \autoref{sec:dhcp-psscope}). They can also be used when you simply prefer PowerShell.}

%--------------------------------------------------------------
\subsection{Use of PowerShell in the Event of DHCP Console Failure}
\label{sec:dhcp-ps}
A failing DHCP console does not mean that new DHCP scopes cannot be configured or removed. PowerShell commands allow DHCP scopes to be configured and modified with a similar degree of control as the console. Some of the important commands for configuring DHCP IPv4 scopes are listed in \autoref{tab:dhcp-cmdlets} \cite{msft-dhcpserver-module}, and examples of their output are shown in \autoref{fig:dhcp-19} to \autoref{fig:dhcp-22}.

\begin{table}[H]
  \centering
  \small\sffamily
  \renewcommand{\arraystretch}{1.25}
  \begin{tabular}{@{}p{7.2cm}p{7cm}@{}}
    \toprule
    \textbf{Command} & \textbf{Purpose} \\
    \midrule
    \texttt{\footnotesize Add-DhcpServerv4Scope} & Add scopes \\
    \texttt{\footnotesize Get-DhcpServerv4Scope} & View scopes \\
    \texttt{\footnotesize Set-DhcpServerv4Scope} & Modify scopes \\
    \texttt{\footnotesize Remove-DhcpServerv4Scope} & Delete scopes \\
    \texttt{\footnotesize Get-DhcpServerv4ScopeStatistics} & Show usage statistics \\
    \texttt{\footnotesize Add-DhcpServerv4Lease} & Add a lease \\
    \texttt{\footnotesize Get-DhcpServerv4Lease} & View leases \\
    \texttt{\footnotesize Remove-DhcpServerv4Lease} & Remove a lease \\
    \texttt{\footnotesize Add-DhcpServerv4Reservation} & Add reservations \\
    \texttt{\footnotesize Get-DhcpServerv4Reservation} & View reservations \\
    \texttt{\footnotesize Set-DhcpServerv4Reservation} & Modify reservations \\
    \texttt{\footnotesize Remove-DhcpServerv4Reservation} & Delete a reservation \\
    \bottomrule
  \end{tabular}
  \caption{PowerShell commands for managing DHCP IPv4 scopes}\label{tab:dhcp-cmdlets}
\end{table}

\shot{dhcp_19.png}{Example of Accessing DHCP Through PowerShell}{fig:dhcp-19}
\shot{dhcp_20.png}{Example of Adding a DHCP Reservation}{fig:dhcp-20}
\shot{dhcp_21.png}{Example of Verification of a DHCP Reservation}{fig:dhcp-21}
\shot{dhcp_22.png}{Example of IP Address Tracking}{fig:dhcp-22}

%--------------------------------------------------------------
\subsection{Using Windows PowerShell to Create a DHCP Scope}
\label{sec:dhcp-psscope}
To create a DHCP scope using Windows PowerShell, it must be opened as Administrator \cite{sherry-dhcp-scope}. Considering that scopes have most likely been made already, a command must first be run to check for existing scopes.

The command \texttt{Get-DhcpServerv4Scope} lists every scope on the DHCP server that was created using the DHCP console (\autoref{fig:dhcp-23}) \cite{sherry-dhcp-scope}. If there are no preexisting scopes, nothing is shown.

\shot{dhcp_23.png}{Example of the Get-DhcpServerv4Scope Command in Use}{fig:dhcp-23}

Creating a scope in PowerShell would normally involve a lot of typing, as the \texttt{Add-DhcpServerv4Scope} command has many parameters. A conventional use of the full command would look like the following:

\begin{lstlisting}
Add-DhcpServerv4Scope `
    -Name "Test Scope" `
    -Description "Scope for DHCP testing" `
    -StartRange 192.168.10.100 `
    -EndRange 192.168.10.200 `
    -SubnetMask 255.255.255.0 `
    -LeaseDuration 0.08:00:00 `
    -State Active `
    -PassThru
\end{lstlisting}

This can be tedious, so splatting can be used to make passing parameters easier \cite{listek-splatting}. PowerShell splatting is when a hashtable is used to create a parameter set that is passed to the cmdlet, which then creates the DHCP scope according to the given parameters. For example, if the following were entered in the PowerShell console:

\begin{lstlisting}
$HashArgs = @{
    'Name'          = 'Head Office Wireless Scope'
    # This sets the name of the wireless scope
    'Description'   = 'Wireless Clients'
    # This is the human-readable description of the scope
    'StartRange'    = '10.1.0.100'
    # Specifies the starting IP address in the scope
    'EndRange'      = '10.1.0.200'
    # Specifies the end IP address in the scope
    'SubnetMask'    = '255.255.255.0'
    # Specifies the subnet mask of the scope
    'State'         = 'Active'
    # Activates the scope
    'LeaseDuration' = '1.00:00:00'
    # Specifies the length of the lease duration
}
Add-DhcpServerv4Scope @HashArgs
\end{lstlisting}

the hashtable \texttt{\$HashArgs} is given the DHCP scope parameters. When the \texttt{Add-\allowbreak DhcpServerv4Scope} command is entered, a new scope named Head Office Wireless Scope is made, with an address pool between 10.1.0.100 and 10.1.0.200 and a subnet mask of 255.255.255.0 \cite{sherry-dhcp-scope}. The lease duration is one day and the scope is made active. To check that the scope was made, run \texttt{Get-DhcpServerv4Scope} again. A way to make the scopes easier to read is to add \texttt{| Format-List} at the end of the command \cite{sherry-dhcp-scope}.

%--------------------------------------------------------------
\section{Testing}
\label{sec:dhcp-testing}

\subsection{Checking the Server}
\label{sec:dhcp-checkserver}
To test the operation of a DHCP server, there must be confirmation that data packets are travelling between the server and client devices. To start, the server is checked in a PowerShell window that is run as Administrator, using the command:

\begin{lstlisting}
Get-Service -Name DHCPServer
\end{lstlisting}

The Status column must display \emph{Running}. If the command shows the service as \emph{Stopped}, the server cannot answer any client requests. The reason may be that the service has not started or that the server has not been authorised. If it needs to be started, run

\begin{lstlisting}
Start-Service -Name DHCPServer
\end{lstlisting}

and then run the \texttt{Get-Service} command again to confirm that it is running. Another probable reason is that the server needs authorisation. The command \texttt{Get-DhcpServerInDC} verifies whether the DHCP server is authorised in Active Directory \cite{msft-dhcp-quickstart}. For a server that is not yet listed, it must first be added to the directory with:

\begin{lstlisting}
Add-DhcpServerInDC -DnsName <server FQDN> -IPAddress <server IP address>
\end{lstlisting}

A few other checks are to confirm the configuration and the scope usage, seeing how many addresses are in use and how many are free \cite{msft-dhcp-quickstart}:

\begin{lstlisting}
Get-DhcpServerv4ScopeStatistics
Get-DhcpServerv4Scope
\end{lstlisting}

\subsection{Testing from a Client and Verifying on the Server}
\label{sec:dhcp-testclient}
To test the server from a client, the client has to request a lease. On a Windows client, open Command Prompt and run \texttt{ipconfig /release}, then \texttt{ipconfig /renew}, then \texttt{ipconfig /all}. The first command releases the DHCP client lease, the second renews it, and \texttt{ipconfig /all} (or just \texttt{ipconfig}) verifies the DHCP client lease (\autoref{fig:dhcp-24}) \cite{msft-dhcp-renew}.

\shot{dhcp_24.png}{\texttt{ipconfig /all} Command in Command Prompt}{fig:dhcp-24}

The address, subnet mask, gateway, DNS servers and DHCP server line must all match the scope (\autoref{tab:dhcp-config}).

\issue{Client Has Incorrect or Missing Options}
{The client lease shows incorrect or missing options. This can happen because the network configuration is not working correctly, because the client cannot support the missing options, or because the server is not configured to hand them out \cite{msft-dhcp-clients}.}
{Change the list of values for the missing or incorrect options at the DHCP scope in use. You might have to configure the options as a Server Option at the DHCP server, remove it, and set the correct value in the Scope Options of the DHCP scope that is used by the client \cite{msft-dhcp-clients}.}

After the lease is recorded, run the following command to confirm that clients are receiving addresses. The cmdlet also has a means of identifying bad leases with \texttt{-BadLeases}, and \texttt{-AllLeases} specifies all IPv4 address leases \cite{msft-get-lease}:

\begin{lstlisting}
Get-DhcpServerv4Lease -ScopeId <scope ID>
\end{lstlisting}

The server logs and lease records are kept in two places: major DHCP service events are tracked in Event Viewer, but individual lease events go to separate DHCP log files. By default, these are in the directory \path{C:\Windows\System32\dhcp}, with one file per weekday \cite{garn-dhcp-logs}.

%==============================================================
%  CHAPTER 5 - DNS      (skeleton: fill in and delete the TODO boxes)
%  Pattern: copy any \shot line, change the file name, caption and label.
%  Screenshots go in images/ using the file name written here (dns_01.png, ...).
%==============================================================
\chapter{DNS Configuration}
\label{chap:dns}

\todo{Chapter 5 was added late, so part of the DNS work was already documented in Chapter~\ref{chap:adds}, and \textbf{does not need to be redone here}. Refer to it instead of repeating it: \autoref{sec:step3} (Chapter~\ref{chap:adds}) covers opening DNS Manager, checking the forward lookup zone and creating the reverse lookup zone (see \autoref{fig:dnsopen} and \autoref{fig:zonescope}), and \autoref{sec:step4} covers the Host (A) and PTR records, the adapter DNS setting and the \texttt{nslookup} tests (see \autoref{fig:arecord}, \autoref{fig:adapterdns} and \autoref{fig:nslookupbasic}); the values are in \autoref{tab:config}. In this chapter, write the DNS background (what it is, why it matters, the risks), summarise those steps in a few sentences with a pointer to Chapter~\ref{chap:adds}, and only add screenshots and tests for DNS work that is \emph{not} already shown there. Delete this box when done.}

\section{What is DNS?}
\label{sec:dns-what}
\todo{Write 3--5 sentences defining the Domain Name System (DNS), zones and records, and cite a source with \texttt{\textbackslash cite\{key\}}.}

\section{Why is DNS Important?}
\label{sec:dns-why}
\todo{Explain why the organisation needs it (name resolution, locating domain controllers) and how it fits with AD DS and DHCP.}

\section{Security Risks if Not Set Up or Set Up Incorrectly}
\label{sec:dns-risks}
\todo{Describe the risks of leaving the role out or misconfiguring it (e.g.\ failed logons, insecure dynamic updates).}

\section{Step-by-Step Configuration Guide}
\label{sec:dns-guide}
\todo{Add a short table of the key values used (see Table~\ref{tab:config} in Chapter~\ref{chap:adds} for the format), then one subsection per step. Each step: one paragraph of text that refers to its screenshot with \texttt{\textbackslash autoref\{fig:dns-01\}}, followed by the \texttt{\textbackslash shot} line.}

\subsection{Step 1: Verify the DNS Server Role}
\label{sec:dns-step1}
\todo{Describe the step here.}

\shot{dns_01.png}{Replace this caption (DNS, step 1)}{fig:dns-01}
\shot{dns_02.png}{Replace this caption (DNS, step 1 continued)}{fig:dns-02}

\subsection{Step 2: Configure the Zones and Records}
\label{sec:dns-step2}
\todo{Describe the step here.}

\shot{dns_03.png}{Replace this caption (DNS, step 2)}{fig:dns-03}

\section{Testing}
\label{sec:dns-testing}
\todo{Explain how the service was tested and what the result was.}

\shot{dns_test_01.png}{Replace this caption (DNS test result)}{fig:dns-test-01}

\section{Troubleshooting: Possible Sticking Points \& Fixes}
\label{sec:dns-troubleshooting}
\issue{1. Replace with the problem title}
{Describe what the user sees.}
{Describe how to fix it.}

\chapter{Group Policy Setup and Configuration for Windows Server 2019}
\label{chap:gpo}

\section{Introduction}
\label{sec:gpo-intro}
Group Policy is a feature of Microsoft Windows that provides centralised management and configuration of operating systems, applications, and users' settings in an Active Directory environment. It allows administrators to define security settings, software installation policies, and user interface configurations across multiple computers within a domain.

A well-implemented Group Policy can reduce risk, improve security, and ensure consistency across the network. Group Policy Objects (GPOs) can be defined for specific users, groups, or computers, and can be linked to organisational units (OUs) to apply settings based on the structure of Active Directory \cite{msft-gpo-sec}.

This chapter covers the setup and configuration of Group Policy in a Windows Server 2019 environment, including the creation of an OU structure and the configuration of the GPO settings required by the project, followed by testing of their application.

\section{How Does Group Policy Work?}
\label{sec:gpo-how}
Group Policy relies on the Active Directory infrastructure to apply settings to users and computers. When a user logs in or a computer starts up, the Group Policy engine retrieves the relevant GPOs from the domain controller and applies the settings defined within those GPOs. GPOs are processed in the following order:
\begin{enumerate}
  \item \textbf{Local Group Policy:} settings defined on the local computer.
  \item \textbf{Site-level GPOs:} settings defined for the site in which the computer resides.
  \item \textbf{Domain-level GPOs:} settings defined for the entire domain.
  \item \textbf{Organisational Unit (OU)-level GPOs:} settings defined for the specific OU in which the user or computer resides.
\end{enumerate}

When multiple GPOs are applied, the settings from the GPO processed last (the highest precedence) override those from the earlier ones. This allows administrators to create a hierarchy of policies that can be tailored to specific users, groups, or computers.

Policies are refreshed periodically (every 90 minutes by default), and administrators can also force a Group Policy update using the \texttt{gpupdate} command. This ensures that any changes made to GPOs are applied immediately.

\section{Assignment-Specified Group Policy Configuration}
\label{sec:gpo-spec}
This section covers the specific Group Policy configuration required for the assignment. The policies required include disabling the Command Prompt, restricting removable storage access, preventing software installation, and enforcing a minimum password length and a maximum password age. Each of these policies is configured through the Group Policy Management Console (GPMC) and applies to the users and computers placed in the appropriate organisational units (OUs) within Active Directory. \autoref{tab:gpo-policies} summarises the policies configured in this chapter.

\begin{table}[H]
  \centering
  \small\sffamily
  \renewcommand{\arraystretch}{1.3}
  \begin{tabular}{@{}p{5.2cm}p{5.6cm}p{3.3cm}@{}}
    \toprule
    \textbf{Policy} & \textbf{Purpose} & \textbf{Setting used} \\
    \midrule
    Prevent access to the command prompt & Blocks the interactive command prompt (\texttt{cmd.exe}) & Enabled \\
    Maximum password age & Forces users to change their password regularly & 42 days \\
    Minimum password length & Prevents short, easily guessed passwords & 10 characters \\
    All Removable Storage classes: Deny all access & Blocks USB drives and other removable media & Enabled \\
    Turn off Windows Installer & Prevents unauthorised software installation & Enabled (Always) \\
    \bottomrule
  \end{tabular}
  \caption{Group Policy settings configured in this chapter}\label{tab:gpo-policies}
\end{table}

\section{Organisational Unit Creation and Configuration}
\label{sec:gpo-ou}
The workflow for creating and configuring Group Policy Objects (GPOs) involves several steps, including accessing the Group Policy Management Console (GPMC), creating a new GPO, configuring the desired settings, and linking the GPO to the appropriate organisational unit (OU). Therefore, the OU structure must be defined in Active Directory before the GPOs are created and linked. \autoref{tab:gpo-objects} lists the Active Directory objects created for the tests.

\begin{table}[H]
  \centering
  \small\sffamily
  \renewcommand{\arraystretch}{1.3}
  \begin{tabular}{@{}p{3.4cm}p{4.4cm}p{6.3cm}@{}}
    \toprule
    \textbf{Object} & \textbf{Name} & \textbf{Details} \\
    \midrule
    Organisational unit & Group1-OU & Created under OPS.local; protected from accidental deletion \\
    Users & G1 User1, G1 User2 & Logon names of the form G1User1@OPS.local; must change password at next logon \\
    Security group & G1 Users & Global scope, security type; contains both users \\
    \bottomrule
  \end{tabular}
  \caption{Active Directory objects created for the Group Policy tests}\label{tab:gpo-objects}
\end{table}

%--------------------------------------------------------------
\subsection{Creating the Organisational Unit}
\label{sec:gpo-create-ou}
Open Active Directory Users and Computers from Server Manager (\autoref{fig:gpo-01}), right-click the domain OPS.local, and choose New $\rightarrow$ Organizational Unit (\autoref{fig:gpo-02}). Name the new OU \emph{Group1-OU} (\autoref{fig:gpo-03}). The OU is then listed under the selected domain (\autoref{fig:gpo-04}).

\shot{gpo_01.png}{Active Directory Users and Computers}{fig:gpo-01}
\shot{gpo_02.png}{New OU Object Selected Under the Domain OPS.local}{fig:gpo-02}
\shot{gpo_03.png}{New OU Object Created and Named Group1-OU}{fig:gpo-03}
\shot{gpo_04.png}{New OU Object Listed Under the Selected Domain}{fig:gpo-04}

%--------------------------------------------------------------
\subsection{Creating the Users}
\label{sec:gpo-create-users}
Right-click Group1-OU and choose New $\rightarrow$ User (\autoref{fig:gpo-05}). Enter the user's name and logon name (\autoref{fig:gpo-06}), then set an initial password and the account options (\autoref{fig:gpo-07}). A summary of the new object is shown before it is created (\autoref{fig:gpo-08}). The process was repeated for a second user, and both users are listed in the OU (\autoref{fig:gpo-09}).

\shot{gpo_05.png}{New User Object Selected Under the Group1-OU}{fig:gpo-05}
\shot{gpo_06.png}{New User Object Configuration}{fig:gpo-06}
\shot{gpo_07.png}{New User Object Password Configuration}{fig:gpo-07}
\shot{gpo_08.png}{New User Object Summary}{fig:gpo-08}
\shot{gpo_09.png}{User Objects Created and Named; Two Users Are Listed}{fig:gpo-09}

%--------------------------------------------------------------
\subsection{Creating the Security Group}
\label{sec:gpo-create-group}
To allow users to be managed together, a security group was created in the OU. Right-click Group1-OU and choose New $\rightarrow$ Group (\autoref{fig:gpo-10}), then configure the group name, scope, and type (\autoref{fig:gpo-11}). The new group is listed under Group1-OU (\autoref{fig:gpo-12}).

\shot{gpo_10.png}{New Group Object Selected Under the Group1-OU (Allows Grouping of Users)}{fig:gpo-10}
\shot{gpo_11.png}{New Group Object Configuration: Group Name, Scope and Type}{fig:gpo-11}
\shot{gpo_12.png}{New Group Object Listed Under the Group1-OU}{fig:gpo-12}

%--------------------------------------------------------------
\subsection{Adding Users to the Group}
\label{sec:gpo-add-users}
There are two ways to configure group membership. \textbf{Method~1} starts from the user: open the user's Properties (\autoref{fig:gpo-13}), go to the Member Of tab (\autoref{fig:gpo-14} and \autoref{fig:gpo-15}), click Add (\autoref{fig:gpo-16}), search for the group (\autoref{fig:gpo-17} and \autoref{fig:gpo-18}), and confirm the selection (\autoref{fig:gpo-19}). The membership tab then lists the group (\autoref{fig:gpo-20}). This method was used for G1 User1.

\shot{gpo_13.png}{User1 Showing Membership Configuration Method 1}{fig:gpo-13}
\shot{gpo_14.png}{User1 Membership Configuration, Step 1}{fig:gpo-14}
\shot{gpo_15.png}{User1 Membership Configuration, Step 2}{fig:gpo-15}
\shot{gpo_16.png}{User1 Membership Configuration, Step 3}{fig:gpo-16}
\shot{gpo_17.png}{Group Search Window}{fig:gpo-17}
\shot{gpo_18.png}{G1 Users Group Find and Selection}{fig:gpo-18}
\shot{gpo_19.png}{Group Selection and Addition Window}{fig:gpo-19}
\shot{gpo_20.png}{User1 Membership Tab}{fig:gpo-20}

\textbf{Method~2} starts from the user or the group. Right-click the user and choose Add to a group (\autoref{fig:gpo-21}), or open the properties of the G1 Users group to add users to it (\autoref{fig:gpo-22}). On the Members tab (\autoref{fig:gpo-23}), click Add, then find and select the users (\autoref{fig:gpo-24} and \autoref{fig:gpo-25}). The members list now includes User2 (\autoref{fig:gpo-26}). This method was used for G1 User2.

\shot{gpo_21.png}{User2 Showing Membership Configuration Method 2: Add to a Group Selected}{fig:gpo-21}
\shot{gpo_22.png}{G1 Users Group Properties Selected to Add Users to It}{fig:gpo-22}
\shot{gpo_23.png}{Properties Tab With List of Members for the G1 Users Group}{fig:gpo-23}
\shot{gpo_24.png}{Users Find and Selection}{fig:gpo-24}
\shot{gpo_25.png}{User2 Selected and Added to the Group}{fig:gpo-25}
\shot{gpo_26.png}{Properties Tab With List of Members for the G1 Users Group; the List Now Includes User2}{fig:gpo-26}

\section{Group Policy Object Creation and Configuration}
\label{sec:gpo-gpo}
With the OU structure in place, the policies were configured in the Group Policy Management Editor, editing the GPO named \emph{Group1 Domain Policies}. The GPO is created in the Group Policy Management Console and linked to Group1-OU so that its settings apply to the users and computers in that OU.

\subsection{Prevent Access to the Command Prompt}
\label{sec:gpo-cmd}
The policy \emph{Prevent access to the command prompt} was located under User Configuration $\rightarrow$ Administrative Templates $\rightarrow$ System (\autoref{fig:gpo-27}). It was set to Enabled (\autoref{fig:gpo-28}), with the option to also disable command prompt script processing left at No. The policy is shown as enabled in the policy list (\autoref{fig:gpo-29}).

\shot{gpo_27.png}{Group Policy Management Editor: Prevent Access to the Command Prompt Selected}{fig:gpo-27}
\shot{gpo_28.png}{Policy Settings for Prevent Access to the Command Prompt: Enabled Selected}{fig:gpo-28}
\shot{gpo_29.png}{Prevent Access to the Command Prompt Shown as Enabled}{fig:gpo-29}

\subsection{Password Policy}
\label{sec:gpo-password}
The password rules are found under Computer Configuration $\rightarrow$ Windows Settings $\rightarrow$ Security Settings $\rightarrow$ Account Policies $\rightarrow$ Password Policy. \emph{Maximum password age} was selected (\autoref{fig:gpo-30}) and set to 42 days (\autoref{fig:gpo-31}), which is confirmed in the policy list (\autoref{fig:gpo-32}). \emph{Minimum password length} was set to 10 characters (\autoref{fig:gpo-33}), also confirmed in the policy list (\autoref{fig:gpo-34}).

\shot{gpo_30.png}{Password Policy Settings: Maximum Password Age Selected}{fig:gpo-30}
\shot{gpo_31.png}{Maximum Password Age Set to 42 Days}{fig:gpo-31}
\shot{gpo_32.png}{Maximum Password Age Shown as 42 Days}{fig:gpo-32}
\shot{gpo_33.png}{Minimum Password Length Set to 10 Characters}{fig:gpo-33}
\shot{gpo_34.png}{Minimum Password Length Shown as 10 Characters}{fig:gpo-34}

\subsection{Removable Storage Access}
\label{sec:gpo-removable}
Under Administrative Templates $\rightarrow$ System $\rightarrow$ Removable Storage Access, the policy \emph{All Removable Storage classes: Deny all access} was selected (\autoref{fig:gpo-35}) and set to Enabled (\autoref{fig:gpo-36}). When it is enabled, no access is allowed to any removable storage class. The policy is shown as enabled in the policy list (\autoref{fig:gpo-37}).

\shot{gpo_35.png}{All Removable Storage Classes: Deny All Access Selected}{fig:gpo-35}
\shot{gpo_36.png}{All Removable Storage Classes: Deny All Access Enabled}{fig:gpo-36}
\shot{gpo_37.png}{All Removable Storage Classes: Deny All Access Shown as Enabled}{fig:gpo-37}

\subsection{Preventing Software Installation}
\label{sec:gpo-installer}
Under Administrative Templates $\rightarrow$ Windows Components $\rightarrow$ Windows Installer, the policy \emph{Turn off Windows Installer} was selected (\autoref{fig:gpo-38}) and set to Enabled with the option \emph{Always} (\autoref{fig:gpo-39}), which stops users from installing software with Windows Installer packages. The policy is shown as enabled in the policy list (\autoref{fig:gpo-40}).

\shot{gpo_38.png}{Turn off Windows Installer Selected}{fig:gpo-38}
\shot{gpo_39.png}{Turn off Windows Installer Enabled}{fig:gpo-39}
\shot{gpo_40.png}{Turn off Windows Installer Shown as Enabled}{fig:gpo-40}

\section{Testing}
\label{sec:gpo-testing}
Testing is a critical step in the Group Policy configuration process to ensure that the policies are applied correctly and function as intended. Three methods are available:
\begin{enumerate}
  \item \textbf{Pre-deployment simulation.} After creating and linking the GPOs, administrators can use simulation tools to test the policies before applying them in a production environment.
  \item \textbf{The \texttt{gpresult} command.} Administrators can verify that the policies are being applied to the target users and computers.
  \item \textbf{Logging in as an affected user.} Logging in as a user affected by the GPOs confirms that the restrictions and settings are enforced.
\end{enumerate}

The first and last methods are easily understood and are therefore not covered in detail in this report. The second method is described below. On the client machine, first refresh the policies and then report which GPOs were applied:

\begin{lstlisting}
gpupdate /force
gpresult /r
\end{lstlisting}

The output of \texttt{gpresult /r} lists the applied GPOs under \emph{Applied Group Policy Objects} for both the computer and the user. The restrictions can then be confirmed by logging in as one of the test users and attempting the blocked actions, such as opening Command Prompt or using a removable drive.

\section{Chapter Summary}
\label{sec:gpo-summary}
This chapter described how Group Policy works in an Active Directory environment and how the policies required by the project were configured. An organisational unit (Group1-OU), two test users, and a security group (G1 Users) were created in Active Directory, and policies were configured to disable the Command Prompt, deny access to all removable storage, turn off Windows Installer, enforce a minimum password length of 10 characters, and require passwords to be changed every 42 days. The testing procedure to verify that the policies are applied was also described.

%==============================================================
%  CHAPTER 7 - REMOTE ACCESS      (skeleton: fill in and delete the TODO boxes)
%  Pattern: copy any \shot line, change the file name, caption and label.
%  Screenshots go in images/ using the file name written here (remote_01.png, ...).
%==============================================================
\chapter{Remote Access Configuration}
\label{chap:remote}

\section{What is Remote Access?}
\label{sec:remote-what}
\todo{Write 3--5 sentences defining the Remote Access role (and VPN / DirectAccess, whichever was configured) and cite a source with \texttt{\textbackslash cite\{key\}}.}

\section{Why is Remote Access Important?}
\label{sec:remote-why}
\todo{Explain why the organisation needs it (secure connectivity for remote users) and how it fits with AD DS and Group Policy.}

\section{Security Risks if Not Set Up or Set Up Incorrectly}
\label{sec:remote-risks}
\todo{Describe the risks of leaving the role out or misconfiguring it.}

\section{Step-by-Step Configuration Guide}
\label{sec:remote-guide}
\todo{Add a short table of the key values used (see Table~\ref{tab:config} in Chapter~\ref{chap:adds} for the format), then one subsection per step. Each step: one paragraph of text that refers to its screenshot with \texttt{\textbackslash autoref\{fig:remote-01\}}, followed by the \texttt{\textbackslash shot} line.}

\subsection{Step 1: Install the Remote Access Role}
\label{sec:remote-step1}
\todo{Describe the step here.}

\shot{remote_01.png}{Replace this caption (Remote Access, step 1)}{fig:remote-01}
\shot{remote_02.png}{Replace this caption (Remote Access, step 1 continued)}{fig:remote-02}

\subsection{Step 2: Configure the Service}
\label{sec:remote-step2}
\todo{Describe the step here.}

\shot{remote_03.png}{Replace this caption (Remote Access, step 2)}{fig:remote-03}

\section{Testing}
\label{sec:remote-testing}
\todo{Explain how the service was tested and what the result was.}

\shot{remote_test_01.png}{Replace this caption (Remote Access test result)}{fig:remote-test-01}

\section{Troubleshooting: Possible Sticking Points \& Fixes}
\label{sec:remote-troubleshooting}
\issue{1. Replace with the problem title}
{Describe what the user sees.}
{Describe how to fix it.}

%==============================================================
%  CHAPTER 8 - WINDOWS CERTIFICATE SERVICES
%  Screenshots: images/certs_NN.png  =  WinCertServNN.png from the supplied zip
%  (WinCertServ17 duplicates WinCertServ16 and is not used;
%   WinCertServ22 is used as images/certs_test_01.png).
%  New bibliography entries are listed in the cover note -- add them to the
%  bibliography before compiling, or the \cite keys below will show as [?].
%==============================================================
\chapter{Windows Certificate Services Configuration}
\label{chap:certs}

\section{What are Windows Certificate Services?}
\label{sec:certs-what}
Active Directory Certificate Services (AD CS) is a Windows Server role that is used to issue and manage the public key infrastructure (PKI) certificates that secure communication and authentication on a network \cite{certs-overview}. At the centre of AD CS is the certification authority (CA). A CA accepts certificate requests, verifies the identity of the requester according to its policy, and then signs the certificate with its own private key. It is also responsible for revoking certificates that are no longer valid and for publishing revocation lists \cite{certs-cc731523}. Because every certificate it signs can be traced back to the CA, clients that trust the CA will also trust the certificates it issues.

AD CS is made up of several role services. Besides the CA itself, this project installed Certification Authority Web Enrollment, an Online Responder, the Certificate Enrollment Policy Web Service and the Certificate Enrollment Web Service. The Online Responder answers certificate revocation checks using the Online Certificate Status Protocol (OCSP), and the two enrollment web services allow certificates to be requested over HTTPS, including by computers that are not connected to the domain \cite{certs-cc731523,certs-overview}.

\section{Why are Certificate Services Important?}
\label{sec:certs-why}
For an organisation such as CPUT, certificates provide the trust that other services depend on. They allow users, computers and services to prove who they are, and they make it possible to encrypt traffic so that it cannot be read in transit. Typical uses include secure wireless networks, virtual private networks (VPNs), IPsec, the Encrypting File System and smart card logon \cite{certs-cc731523}.

AD CS also fits closely with the other services built in this project. An enterprise CA relies on Active Directory Domain Services (Chapter~\ref{chap:adds}): identity information that is already stored in Active Directory is used when certificates are issued, and Group Policy can be used to decide which users and computers may receive which types of certificate \cite{certs-overview}. DNS is needed so that clients can locate the CA and its web services by name, and the certificates issued by the CA can in turn be used to protect the Remote Access service configured earlier in this report. Running one internal CA means that the organisation does not have to buy certificates from an external provider for every internal server.

\section{Security Risks if Not Set Up or Set Up Incorrectly}
\label{sec:certs-risks}
If no CA is deployed, internal services have no trusted way of proving their identity. Administrators are then forced to use unencrypted connections, or self-signed certificates that every user must be told to ignore warnings for, which teaches users to accept certificate errors and makes impersonation of a server much easier.

If the CA is set up incorrectly, the damage is greater, because everything that trusts the CA inherits its weaknesses. A short key length or an outdated hash algorithm (the wizard still lists SHA1 and MD5, see Figure~\ref{fig:certs-12}) makes certificates easier to forge \cite{certs-docs}. A CA whose private key is stolen or whose server is compromised can be used to issue valid certificates to an attacker. A CA validity period that is shorter than the certificates it issues, or revocation information that is unavailable, means that certificates fail unexpectedly or that revoked certificates are still accepted. Finally, the name and domain membership of a server cannot be changed once it hosts a CA (Figure~\ref{fig:certs-03}), so a mistake made before installation is expensive to correct afterwards. Web enrollment that is left without a valid HTTPS certificate would send credentials over an unprotected channel.

\section{Step-by-Step Configuration Guide}
\label{sec:certs-guide}
The configuration is carried out in three steps: installing the role (subsection~\ref{sec:certs-step1}), configuring the certification authority and its web services (subsection~\ref{sec:certs-step2}), and configuring the Certificate Enrollment Web Service in a second pass of the wizard (subsection~\ref{sec:certs-step3}). Table~\ref{tab:certs-config} lists the values used throughout this guide. If anything goes wrong, see section~\ref{sec:certs-troubleshooting}.

\begin{table}[htbp]
  \centering
  \small
  \begin{tabular}{@{}ll@{}}
    \toprule
    \textbf{Setting} & \textbf{Value used} \\
    \midrule
    Server name                        & Group1-OPS-SERVER (Group1-OPS-SERVER.OPS.local) \\
    Domain                             & OPS.local \\
    Account used for configuration     & OPS\textbackslash Administrator \\
    CA setup type                      & Enterprise CA \\
    CA type                            & Root CA \\
    Private key                        & New private key \\
    Cryptographic provider             & RSA\#Microsoft Software Key Storage Provider \\
    Key length / hash algorithm        & 2048 bits / SHA256 \\
    CA common name                     & OPS-GROUP1-OPS-SERV-CA \\
    Distinguished name suffix          & DC=OPS,DC=local \\
    CA validity period                 & 5 years (expires 2031/09/08) \\
    Database and log location          & C:\textbackslash Windows\textbackslash system32\textbackslash CertLog \\
    Authentication (CEP and CES)       & Windows integrated authentication \\
    CES service account                & Built-in application pool identity \\
    \bottomrule
  \end{tabular}
  \caption{Key configuration values used in this chapter}
  \label{tab:certs-config}
\end{table}

\subsection{Step 1: Install the Role}
\label{sec:certs-step1}
Open Server Manager, choose \emph{Manage} $\rightarrow$ \emph{Add Roles and Features}, and continue to the \emph{Select server roles} page. Tick \emph{Active Directory Certificate Services}; the roles installed in earlier chapters (AD DS, DHCP and DNS) are already marked as installed (\autoref{fig:certs-01}). No additional features are required, so the \emph{Select features} page is accepted unchanged (\autoref{fig:certs-02}).

\shot{certs_01.png}{Selecting the Active Directory Certificate Services Role}{fig:certs-01}
\shot{certs_02.png}{Select Features Page Left at Its Defaults}{fig:certs-02}

The introduction page for AD CS warns that the computer name and domain settings cannot be changed after a certification authority has been installed (\autoref{fig:certs-03}). For this reason the server was already named and joined to the OPS.local domain before this step was started. On the \emph{Select role services} page the Certification Authority, Certification Authority Web Enrollment, Online Responder, Certificate Enrollment Policy Web Service and Certificate Enrollment Web Service were ticked, while the Network Device Enrollment Service was left unticked because no network devices need certificates in this project (\autoref{fig:certs-04}).

\shot{certs_03.png}{Introduction to AD CS and the Note on Computer Name Changes}{fig:certs-03}
\shot{certs_04.png}{Selecting the AD CS Role Services}{fig:certs-04}

After the installation is confirmed, the progress page reports that the installation succeeded but that configuration is still required (\autoref{fig:certs-05}). A warning flag also appears in the Server Manager notifications with a link to configure AD CS on the destination server, which starts the configuration wizard used in the next step (\autoref{fig:certs-06}).

\shot{certs_05.png}{Installation Progress Showing That Further Configuration Is Required}{fig:certs-05}
\shot{certs_06.png}{Post-Deployment Configuration Notification for AD CS}{fig:certs-06}

\subsection{Step 2: Configure the Certification Authority}
\label{sec:certs-step2}
The AD CS Configuration wizard first asks for credentials. An enterprise CA can only be installed by a member of the Enterprise Admins group, so the OPS\textbackslash Administrator account was used (\autoref{fig:certs-07}). On the next page the role services to configure were selected: the Certification Authority, Certification Authority Web Enrollment, Online Responder and Certificate Enrollment Policy Web Service. The Certificate Enrollment Web Service is not selected in this first pass and is configured separately in Step~3 (\autoref{fig:certs-08}).

\shot{certs_07.png}{AD CS Configuration: Credentials}{fig:certs-07}
\shot{certs_08.png}{AD CS Configuration: Role Services to Configure}{fig:certs-08}

The setup type was set to \emph{Enterprise CA}, because the server is a domain member and the CA must integrate with Active Directory (\autoref{fig:certs-09}). It was configured as a \emph{Root CA}, since it is the first and only CA in the hierarchy (\autoref{fig:certs-10}).

\shot{certs_09.png}{Setup Type: Enterprise CA Selected}{fig:certs-09}
\shot{certs_10.png}{CA Type: Root CA Selected}{fig:certs-10}

A CA must have a private key in order to sign certificates, so the option to create a new private key was kept (\autoref{fig:certs-11}). On the cryptography page the default RSA\#Microsoft Software Key Storage Provider was kept with a key length of 2048 bits, and SHA256 was chosen as the hash algorithm used to sign issued certificates. The older SHA1 and MD5 algorithms are shown in the list but were not selected (\autoref{fig:certs-12}).

\shot{certs_11.png}{Private Key: Create a New Private Key}{fig:certs-11}
\shot{certs_12.png}{Cryptography for CA: RSA, 2048-bit Key and SHA256}{fig:certs-12}

The CA was given the common name OPS-GROUP1-OPS-SERV-CA. The distinguished name suffix DC=OPS,DC=local is generated automatically from the domain, giving the full name CN=OPS-GROUP1-OPS-SERV-CA,DC=OPS,DC=local (\autoref{fig:certs-13}). The validity period of the CA certificate was set to five years, which gives an expiry date of 2031/09/08. The wizard notes that this period should be longer than the validity of any certificate the CA issues (\autoref{fig:certs-14}). The certificate database and its log were left in the default location, C:\textbackslash Windows\textbackslash system32\textbackslash CertLog (\autoref{fig:certs-15}).

\shot{certs_13.png}{CA Name: Common Name and Distinguished Name}{fig:certs-13}
\shot{certs_14.png}{Validity Period Set to Five Years}{fig:certs-14}
\shot{certs_15.png}{CA Database and Log Locations}{fig:certs-15}

The Certificate Enrollment Policy Web Service (CEP) needs to know how clients will authenticate to it. Windows integrated authentication was selected, which lets domain computers and users authenticate with their existing credentials (\autoref{fig:certs-16}). On the \emph{Server Certificate} page, \emph{Choose and assign a certificate for SSL later} was selected, because no certificate could be issued until the CA itself existed. The wizard warns that the service will only function once a valid certificate has been assigned (\autoref{fig:certs-18}); this is dealt with in Step~3 and in section~\ref{sec:certs-troubleshooting}.

\shot{certs_16.png}{Authentication Type for the Certificate Enrollment Policy Web Service}{fig:certs-16}
\shot{certs_18.png}{Server Authentication Certificate to Be Assigned Later}{fig:certs-18}

The confirmation page summarises the settings before anything is changed. It lists the CA as an Enterprise Root CA with the RSA key provider, SHA256 hash, 2048-bit key, the validity date and the database locations (\autoref{fig:certs-19}). Scrolling down shows the remaining role services, including the Certificate Enrollment Policy Web Service with Windows Integrated Authentication and the server certificate to be chosen later (\autoref{fig:certs-20}). Clicking \emph{Configure} starts the configuration, and its progress is shown in \autoref{fig:certs-21}.

\shot{certs_19.png}{Confirmation Page: Certification Authority Settings}{fig:certs-19}
\shot{certs_20.png}{Confirmation Page: Remaining Role Services}{fig:certs-20}
\shot{certs_21.png}{Configuration in Progress}{fig:certs-21}

\subsection{Step 3: Configure the Certificate Enrollment Web Service}
\label{sec:certs-step3}
The Certificate Enrollment Web Service (CES) requires an existing CA to send requests to, so it was configured in a second run of the AD CS Configuration wizard once the CA was in place. The wizard opens with an empty credentials box and the \emph{Next} button disabled (\autoref{fig:certs-23}). After entering OPS\textbackslash Administrator, the wizard continues; its steps now include \emph{CA for CES} and \emph{Service Account for CES} (\autoref{fig:certs-24}). Only the Certificate Enrollment Web Service remains to be configured on the role services page; the services configured earlier are shown ticked but are no longer selectable (\autoref{fig:certs-25}).

\shot{certs_23.png}{Second Wizard Run: Credentials Required}{fig:certs-23}
\shot{certs_24.png}{Second Wizard Run: Credentials Entered}{fig:certs-24}
\shot{certs_25.png}{Selecting the Certificate Enrollment Web Service}{fig:certs-25}

On the \emph{CA for CES} page the CA created in Step~2 was selected by name, Group1-OPS-SERVER.OPS.local\textbackslash OPS-GROUP1-OPS-SERV-CA, and renewal-only mode was left off (\autoref{fig:certs-26}). Windows integrated authentication was used again (\autoref{fig:certs-27}), and the service was set to run under the built-in application pool identity rather than a separate service account (\autoref{fig:certs-28}).

\shot{certs_26.png}{CA for CES: Selecting the Target CA}{fig:certs-26}
\shot{certs_27.png}{Authentication Type for the Certificate Enrollment Web Service}{fig:certs-27}
\shot{certs_28.png}{Service Account for CES: Built-in Application Pool Identity}{fig:certs-28}

This time the \emph{Server Certificate} page offers a certificate issued by the new CA to the server itself (Issued To and Issued By both show OPS-GROUP1-OPS-SERV-CA, expiring 2031/09/08). It was selected so that the web service can use SSL (\autoref{fig:certs-29}). The confirmation page lists the target CA, Windows Integrated Authentication, the application pool identity and the chosen server certificate (\autoref{fig:certs-30}), and the configuration is then started (\autoref{fig:certs-31}).

\shot{certs_29.png}{Server Certificate: Certificate Issued by the New CA Selected}{fig:certs-29}
\shot{certs_30.png}{Confirmation Page for the Certificate Enrollment Web Service}{fig:certs-30}
\shot{certs_31.png}{Certificate Enrollment Web Service Configuration in Progress}{fig:certs-31}

\section{Testing}
\label{sec:certs-testing}
% NOTE: no post-installation check (e.g. certsrv.msc, https://localhost/certsrv, or the
% Online Responder snap-in) was supplied. If one is captured, add it here as
% certs_test_02.png using the same \shot pattern.
The configuration wizard reports the result of every role service it configures, so the Results page of the first run was used to verify the installation. As shown in \autoref{fig:certs-test-01}, the Certification Authority, Certification Authority Web Enrollment, Online Responder and Certificate Enrollment Policy Web Service all show \emph{Configuration succeeded}. The page also carries two reminders for the Certificate Enrollment Policy Web Service: a server authentication certificate must be configured before clients can use it, and Group Policy must be applied to direct client enrollment requests to the web service. The first of these was addressed by selecting a certificate issued by the new CA when the Certificate Enrollment Web Service was configured (\autoref{fig:certs-29}).

The existence of a certificate issued by OPS-GROUP1-OPS-SERV-CA in the server certificate list of Step~3 is itself a useful test of the CA, because that certificate could only have been issued by the CA once it was running and trusted by the server (\autoref{fig:certs-29}).

\shot{certs_test_01.png}{Results Page: Role Services Configured Successfully}{fig:certs-test-01}

\section{Troubleshooting: Possible Sticking Points \& Fixes}
\label{sec:certs-troubleshooting}
\issue{1. Credentials Page Will Not Let You Continue}
{The \emph{Next} button on the Credentials page is disabled and the credentials box is empty (\autoref{fig:certs-23}). The page lists which role services need local Administrator rights and which need Enterprise Admins membership.}
{Click \emph{Change\ldots} and enter an account that is a member of the Enterprise Admins group, for example OPS\textbackslash Administrator, as in \autoref{fig:certs-24}. Without it the Enterprise CA and the web services cannot be configured.}

\issue{2. Server Name or Domain Needs to Change After the CA Is Installed}
{The introduction to AD CS states that the computer name and domain settings cannot be changed once a certification authority has been installed (\autoref{fig:certs-03}). A server renamed afterwards would no longer match the names recorded in its certificates.}
{Decide on the final server name and join the server to the domain \emph{before} installing the role. If it must be renamed afterwards, the CA has to be removed and configured again, so this is best avoided.}

\issue{3. Warning About the Server Authentication Certificate}
{The Server Certificate page warns that, for the role service to function, the server must be configured with a valid certificate (\autoref{fig:certs-18}), and the Results page repeats this (\autoref{fig:certs-test-01}). The certificate list on that page is empty because no certificate exists yet.}
{This is expected during the first run, because the CA that issues the certificate does not exist yet. Select \emph{Choose and assign a certificate for SSL later} and continue. Once the CA is running, either bind a certificate issued by it to the HTTPS binding of the site in IIS Manager, or select the certificate when configuring the next web service, as in \autoref{fig:certs-29}.}

\issue{4. Schannel Warnings in the Server Manager Event List}
{After the CA role is configured, warnings with the source \emph{Schannel} and event ID 36886 start to appear in the Events tile of Server Manager (visible in \autoref{fig:certs-26} and the following figures). They typically mean that SSL/TLS has no usable server certificate on the machine.}
{Assign a valid server authentication certificate to the HTTPS binding as described in the previous item. The warnings should no longer be generated for new connections once a certificate is bound.}

\issue{5. Installation Succeeded but the CA Does Not Work}
{The installation progress page says that the installation succeeded, yet no certificates can be issued and a yellow warning flag remains in Server Manager (\autoref{fig:certs-05}, \autoref{fig:certs-06}).}
{Installing the role only copies the files. Click the notification flag and choose \emph{Configure Active Directory Certificate Services on the destination server} to start the configuration wizard, then complete Step~2.}

\issue{6. A Role Service Was Left Unconfigured}
{After the first run, the Results page lists only four role services (\autoref{fig:certs-test-01}); the Certificate Enrollment Web Service is missing because it was not selected on the role services page (\autoref{fig:certs-08}).}
{Open the AD CS Configuration wizard again from the Server Manager notification, enter the Enterprise Admin credentials, select the remaining role service, and point it to the existing CA on the \emph{CA for CES} page (\autoref{fig:certs-26}).}

\chapter{Conclusion}
\label{chap:conclusion}

\section{Summary}
\label{sec:conclusion-summary}
\todo{Replace this box with 1--2 paragraphs summarising what was built: the Windows Server 2019 installation (Chapter~\ref{chap:install}), AD DS (Chapter~\ref{chap:adds}), DHCP (Chapter~\ref{chap:dhcp}), DNS (Chapter~\ref{chap:dns}), Group Policy (Chapter~\ref{chap:gpo}), Remote Access (Chapter~\ref{chap:remote}) and Certificate Services (Chapter~\ref{chap:certs}), and whether each met its requirement.}

\section{Key Findings}
\label{sec:conclusion-findings}
\todo{Replace this box with the main findings: what worked, which problems were met (e.g.\ the DNS delegation warning, the static IP prerequisite, the DHCP console failure) and how they were solved.}

\section{Future Work}
\label{sec:conclusion-future}
\todo{Replace this box with realistic improvements, e.g.\ a second domain controller, DHCP failover, further GPO hardening, or auditing and monitoring.}

\section{Final Remarks}
\label{sec:conclusion-remarks}
\todo{Replace this box with a short closing paragraph.}


%==============================================================
%  BIBLIOGRAPHY  (one shared list - add new entries at the end)
%==============================================================
\newpage
\begin{thebibliography}{99}
\addcontentsline{toc}{chapter}{Bibliography}

% ---- Chapter 3: AD DS ----
\bibitem{msft-adds-overview}
Microsoft Learn. (2023). \textit{Active Directory Domain Services Overview}. Microsoft. Retrieved from \url{https://learn.microsoft.com/en-us/windows-server/identity/ad-ds/get-started/virtual-dc/active-directory-domain-services-overview}

\bibitem{msft-gpo-sec}
Microsoft Learn. (2022). \textit{Group Policy and Security Policy Processing}. Microsoft. Retrieved from \url{https://learn.microsoft.com/en-us/windows-server/identity/ad-ds/manage/group-policy/group-policy-processing}

\bibitem{msft-sec-guidelines}
Microsoft Learn. (2024). \textit{Best Practices for Securing Active Directory}. Microsoft. Retrieved from \url{https://learn.microsoft.com/en-us/windows-server/identity/ad-ds/plan/security-best-practices/best-practices-for-securing-active-directory}

\bibitem{msft-dns-integration}
Microsoft Learn. (2023). \textit{DNS and AD DS Integration}. Microsoft. Retrieved from \url{https://learn.microsoft.com/en-us/windows-server/identity/ad-ds/plan/dns-and-ad-ds}

% ---- Chapter 4: DHCP ----
\bibitem{krause2019}
Krause, J. (2019). \textit{Mastering Windows Server 2019} (2nd ed.). Packt Publishing Ltd.

\bibitem{dhcp-groups}
Microsoft. (2026). \textit{DHCP Groups}. Retrieved from \url{https://learn.microsoft.com/zh-cn/previous-versions/windows/it-pro/windows-server-2003/cc737716(v=ws.10)}

\bibitem{msnoob-dhcp}
Anggara, A. \textit{Install DHCP Server on Windows Server 2016}. Retrieved from \url{https://msnoob.com/install-dhcp-server-on-windows-server-2016.html}

\bibitem{ipme-ipv4}
Bello, N. (2026). \textit{How Many IPv4 Addresses Are There?} Retrieved from \url{https://ipme.co/articles/how-many-ipv4-addresses}

\bibitem{rapidseedbox-lease}
Georgiv, D. (2026). \textit{DHCP Lease Time: How Long Should It Be?} Retrieved from \url{https://www.rapidseedbox.com/blog/dhcp-lease-time#2}

\bibitem{msft-dhcp-dns}
Microsoft. (2026). \textit{DHCP: The IPv6 Addresses of the DNS Server Should Be Reachable}. Retrieved from \url{https://learn.microsoft.com/en-us/previous-versions/windows/it-pro/windows-server-2008-R2-and-2008/ee941121(v=ws.10)?redirectedfrom=MSDN}

\bibitem{msft-wins}
Microsoft. (2026). \textit{Chapter 18 -- Maintaining WINS}. Retrieved from \url{https://learn.microsoft.com/en-us/previous-versions/windows/it-pro/windows-server-2000/bb727108(v=technet.10)}

\bibitem{msft-dhcpserver-module}
Microsoft. (2026). \textit{DhcpServer Module}. Retrieved from \url{https://learn.microsoft.com/en-us/powershell/module/dhcpserver/?view=windowsserver2019-ps}

\bibitem{sherry-dhcp-scope}
Sherry, S. (2021). \textit{Managing Microsoft DHCP Scope and Other Components}. Retrieved from \url{https://adamtheautomator.com/dhcp-scope/}

\bibitem{listek-splatting}
Listek, A. (2020). \textit{PowerShell Splatting: Simplify and Organize Your Scripts}. Retrieved from \url{https://adamtheautomator.com/powershell-splatting/}

\bibitem{msft-dhcp-quickstart}
Microsoft. (2026). \textit{Quickstart: Install and Configure DHCP Server}. Retrieved from \url{https://learn.microsoft.com/en-us/windows-server/networking/technologies/dhcp/quickstart-install-configure-dhcp-server?tabs=powershell}

\bibitem{msft-dhcp-renew}
Microsoft. (2026). \textit{Verify, Release, or Renew a Client Address Lease}. Retrieved from \url{https://learn.microsoft.com/en-us/previous-versions/windows/it-pro/windows-server-2008-R2-and-2008/dd183660(v=ws.10)}

\bibitem{msft-dhcp-clients}
Microsoft. (2026). \textit{Troubleshooting DHCP Clients}. Retrieved from \url{https://learn.microsoft.com/en-us/previous-versions/windows/it-pro/windows-server-2003/cc757164(v=ws.10)?redirectedfrom=MSDN}

\bibitem{msft-get-lease}
Microsoft. (2026). \textit{Get-DhcpServerv4Lease}. Retrieved from \url{https://learn.microsoft.com/fr-fr/previous-versions/windows/powershell-scripting/jj590737(v=wps.620)}

\bibitem{garn-dhcp-logs}
Garn, D. (2024). \textit{Use Advanced DHCP Logging Techniques for Troubleshooting}. Retrieved from \url{https://www.techtarget.com/it-infrastructure/tutorial/Use-advanced-DHCP-logging-techniques-for-troubleshooting}

% ---- Add further references below (Remote Access, Certificate Services, ...) ----

\bibitem{certs-overview}
Microsoft Learn. (2026). \textit{Active Directory Certificate Services Overview}. Microsoft. Retrieved from \url{https://learn.microsoft.com/windows-server/identity/ad-cs/active-directory-certificate-services-overview}

\bibitem{certs-docs}
Microsoft Learn. (2026). \textit{Active Directory Certificate Services Documentation}. Microsoft. Retrieved from \url{https://learn.microsoft.com/windows-server/identity/ad-cs}

\bibitem{certs-cc731523}
Microsoft. (2026). \textit{Active Directory Certificate Services (Windows Server 2008 R2 and 2008)}. Retrieved from \url{https://learn.microsoft.com/en-us/previous-versions/windows/it-pro/windows-server-2008-r2-and-2008/cc731523(v=ws.10)}

\end{thebibliography}

%==============================================================
%  APPENDIX
%==============================================================
\appendix
\chapter{Appendix}
\label{app:appendix}

\section{GitHub Repository}
\todo{Insert the link to the group's GitHub repository here, e.g.\ \texttt{\textbackslash url\{https://github.com/...\}}.}

\end{document}
